\begin{sol}[Cilindro rodando sobre un plano inclinado]

\subsubsection*{Pregunta 1: Aceleración del centro de masa}

Ecuaciones de Newton-Euler para el cilindro, con $x$ a lo largo del plano (positivo hacia abajo) y $C$ el punto de contacto:

\textbf{Traslación} (segunda ley de Newton a lo largo del plano):
$$mg\sin\theta - f = ma_{cm}$$

\textbf{Rotación} respecto al centro de masa $G$, con momento de inercia $I_G = \frac{1}{2}mR^2$ para un cilindro sólido:
$$fR = I_G \alpha = \frac{1}{2}mR^2\alpha$$

\textbf{Restricción de rodadura sin deslizar}:
$$a_{cm} = R\alpha$$

Sustituyendo $\alpha = a_{cm}/R$ en la ecuación de rotación:
$$f = \frac{1}{2}m a_{cm}$$

Sustituyendo en la ecuación de traslación:
$$mg\sin\theta - \frac{1}{2}m a_{cm} = m a_{cm}$$

$$\boxed{a_{cm} = \frac{2}{3}g\sin\theta}$$

\subsubsection*{Pregunta 2: Fuerza de roce estático}

$$f = \frac{1}{2}m a_{cm} = \frac{1}{2}m\left(\frac{2}{3}g\sin\theta\right)$$

$$\boxed{f = \frac{1}{3}mg\sin\theta}$$

\subsubsection*{Pregunta 3: Coeficiente mínimo de roce}

La condición de no deslizamiento requiere $f \leq \mu_s N$. En la dirección perpendicular al plano, el equilibrio da $N = mg\cos\theta$. Entonces:

$$\frac{1}{3}mg\sin\theta \leq \mu_s \, mg\cos\theta$$

$$\boxed{\mu_{s,\text{min}} = \frac{1}{3}\tan\theta}$$

\paragraph{Interpretación física} La aceleración $\frac{2}{3}g\sin\theta$ es menor que $g\sin\theta$ (caso sin rotación, como un bloque deslizando sin roce) porque parte de la energía potencial se convierte en energía cinética rotacional además de traslacional. El resultado es independiente de $m$ y $R$: todos los cilindros sólidos homogéneos ruedan con la misma aceleración en un plano dado.

\end{sol}
